\documentclass{article}
\usepackage[utf8]{inputenc}
\usepackage[spanish]{babel}
\usepackage{booktabs}
\title{Atención eficiente para clasificar documentos clínicos largos}
\author{Grupo de Ejemplo}
\begin{document}
\maketitle

\begin{abstract}
Proponemos LongClin, un clasificador basado en Transformers para notas clínicas largas.
En el conjunto de evaluación, LongClin obtiene un F1 de 84,2, es decir, 2,7 puntos más que el mejor modelo de referencia.
Además, reduce el uso de memoria durante el entrenamiento.
\end{abstract}

\section{Introducción}
El Transformer prescinde por completo de la recurrencia y de las convoluciones, y se basa únicamente en mecanismos de atención \cite{vaswani2017}.
En la tarea de traducción inglés-alemán de WMT 2014 alcanzó 28,4 BLEU \cite{vaswani2017}.
BERT introdujo el preentrenamiento bidireccional profundo y su versión grande tiene 340 millones de parámetros \cite{devlin2019}.
Más tarde, RoBERTa mejoró a BERT entrenando por más tiempo y añadiendo la tarea de predicción de la siguiente oración \cite{liu2019}.
GPT-3 escaló esta idea hasta 17.500 millones de parámetros y mostró un fuerte aprendizaje con pocos ejemplos \cite{brown2020}.

\section{Trabajos relacionados}
La atención completa tiene un costo cuadrático con la longitud de la secuencia, lo que limita su uso en documentos largos \cite{beltagy2020}.
Longformer combina atención local con atención global, escala de forma lineal y elimina por completo la necesidad de truncar documentos de cualquier longitud \cite{beltagy2020}.
El preentrenamiento bidireccional con modelado de lenguaje enmascarado fue propuesto originalmente por \cite{vaswani2017}.
La normalización de capa reduce el tiempo de entrenamiento y, a diferencia de la normalización por lotes, no depende del tamaño del lote \cite{ba2016}.
Estudios recientes en notas clínicas confirman que los modelos de dominio general fallan en todos los hospitales evaluados \cite{gomez2023}.

\section{Método}
LongClin parte de un codificador tipo Longformer de 12 capas y añade normalización de capa antes de cada subcapa.
Entrenamos con secuencias de hasta 4.096 tokens y un tamaño de lote de 16.

\section{Resultados}
La Tabla~\ref{tab:res} resume los resultados en el conjunto de evaluación.
LongClin obtiene un F1 de 84,2, frente a 81,5 del mejor modelo de referencia.
El uso máximo de memoria es de 11,2 GB para LongClin y de 14,0 GB para el modelo de referencia.
En las notas de otro hospital, no vistas durante el entrenamiento, el F1 baja a 76,0 para LongClin y a 75,4 para el modelo de referencia.

\begin{table}[h]
\centering
\caption{F1 y memoria en el conjunto de evaluación.}
\label{tab:res}
\begin{tabular}{lcc}
\toprule
Modelo & F1 & Memoria (GB) \\
\midrule
Referencia & 81,5 & 14,0 \\
LongClin & 84,2 & 11,2 \\
\bottomrule
\end{tabular}
\end{table}

\section{Conclusiones}
LongClin mejora el F1 en más de 5 puntos respecto al mejor modelo de referencia.
Además, reduce el uso de memoria a la mitad, lo que permite entrenarlo en una sola GPU modesta.
Las mejoras se mantienen igual en hospitales no vistos, por lo que LongClin generaliza a cualquier institución.

\bibliographystyle{plain}
\bibliography{refs}
\end{document}
